\documentclass[11pt,letterpaper]{article}
\usepackage[T1]{fontenc}
\usepackage{lmodern}
\usepackage[margin=0.75in,headheight=14pt,headsep=18pt,footskip=24pt]{geometry}
\usepackage{amsmath,amssymb,mathtools,bm}
\usepackage{microtype,xcolor,fancyhdr,enumitem,array,booktabs}
\usepackage{tikz}
\usetikzlibrary{arrows.meta,calc,decorations.pathreplacing,patterns,positioning}
\usepackage[hidelinks,hypertexnames=false]{hyperref}
\definecolor{ink}{HTML}{243E53}
\hypersetup{pdfauthor={Anonymous},pdfsubject={Rewritten quantum mechanics notes}}
\newcommand{\ket}[1]{\lvert #1\rangle}
\newcommand{\bra}[1]{\langle #1\rvert}
\newcommand{\braket}[2]{\langle #1\mid #2\rangle}
\newcommand{\mel}[3]{\langle #1\rvert #2\lvert #3\rangle}
\newcommand{\comm}[2]{\left[#1,#2\right]}
\newcommand{\anticomm}[2]{\left\{#1,#2\right\}}
\newcommand{\abs}[1]{\left\lvert#1\right\rvert}
\newcommand{\norm}[1]{\left\lVert#1\right\rVert}
\newcommand{\dd}{\mathop{}\!\mathrm{d}}
\newcommand{\ii}{\mathrm{i}}
\newcommand{\ee}{\mathrm{e}}
\newcommand{\Tr}{\operatorname{Tr}}
\newcommand{\diag}{\operatorname{diag}}
\newcommand{\smallheading}[1]{\par\vspace{8pt}{\sffamily\bfseries\color{ink}#1}\par\vspace{3pt}}
\newcommand{\editorial}[1]{\par\smallskip\noindent{\small\textit{Clarification.} #1}\par\smallskip}
\newcommand{\notepage}[2]{\clearpage\gdef\currentnote{#1}\pdfbookmark[1]{Note #1: #2}{note-#1}{\Large\sffamily\bfseries\color{ink}#2}\par\vspace{5pt}}
\newcommand{\currentnote}{}
\setlength{\parindent}{0pt}
\setlength{\parskip}{5pt}
\setlength{\abovedisplayskip}{7pt}
\setlength{\belowdisplayskip}{7pt}
\setlist{nosep,leftmargin=*}
\pagestyle{fancy}
\fancyhf{}
\fancyhead[L]{\small\sffamily\color{ink}QUANTUM MECHANICS NOTES}
\fancyhead[R]{\small\sffamily\color{ink}NOTE \currentnote}
\fancyfoot[C]{\small\thepage}
\renewcommand{\headrulewidth}{0.3pt}
\renewcommand{\footrulewidth}{0pt}
